\chapter{Respiratory Pathogens Panel}

\section*{What Is This Test?}
A respiratory pathogens panel (RPP) is a multiplex nucleic acid amplification test (NAAT) that simultaneously detects a defined panel of viral and bacterial respiratory pathogens from a single nasopharyngeal specimen. Panels vary by manufacturer but typically include influenza A and B, respiratory syncytial virus (RSV), parainfluenza viruses 1--4, human metapneumovirus, adenovirus, rhinovirus/enterovirus, coronaviruses (including SARS-CoV-2), \textit{Bordetella pertussis}, \textit{Bordetella parapertussis}, \textit{Chlamydia pneumoniae}, and \textit{Mycoplasma pneumoniae}. The most widely used platforms include BioFire FilmArray Respiratory Panel 2.1 (22 targets, 45-minute run), Luminex NxTAG Respiratory Pathogen Panel (20 targets, 4-hour run), GenMark ePlex Respiratory Pathogen Panel (20+ targets, 2-hour run), and QIAstat-Dx Respiratory Panel (21 targets, 1-hour run). These panels transform the diagnostic approach to respiratory infections by identifying the causative pathogen in 30--70\% of cases where conventional methods (culture, single-target PCR, antigen testing) were negative, enabling targeted therapy and reducing unnecessary antibiotic use.

\section*{Alternative Names}
Multiplex Respiratory PCR; Respiratory Virus Panel; BioFire Respiratory Panel; FilmArray RP; RVP; EPP; Respiratory Pathogen Multiplex; Comprehensive Respiratory PCR

\section*{Principle}
All platforms use multiplex reverse-transcription polymerase chain reaction (RT-PCR) with subsequent detection of amplified nucleic acid. BioFire FilmArray uses nested multiplex PCR with endpoint melt curve analysis in a closed pouch system. Luminex NxTAG uses multiplex PCR followed by bead-based hybridization and fluorescent detection on a flow cytometer. GenMark ePlex uses electrowetting digital microfluidics and competitive DNA hybridization with voltammetric detection. Qiagen QIAstat-Dx uses multiplex PCR with melt curve analysis in a test cartridge. Each platform detects pathogen-specific conserved genetic sequences (e.g., matrix gene for influenza, fusion gene for RSV, hexon gene for adenovirus). Internal controls for specimen adequacy, nucleic acid extraction, and PCR inhibition are included in each cartridge. Results are qualitative (detected/not detected), though semi-quantitative Ct values are displayed on some platforms. Turnaround time ranges from 20--45 minutes (FilmArray) to 4 hours (Luminex).

\section*{History}
The first multiplex respiratory viral panels (xTAG RVP from Luminex) received FDA clearance in 2008. BioFire FilmArray RP received FDA clearance in 2011 and was updated to RP2 in 2017 and RP2.1 in 2020 to include SARS-CoV-2. The addition of SARS-CoV-2 to existing multiplex panels was accelerated by FDA Emergency Use Authorizations in 2020. The evolution from single-plex or low-plex testing to comprehensive multiplex panels has transformed respiratory diagnostics, allowing near-complete syndromic testing in under 1 hour from specimen collection. Clinical studies have demonstrated that multiplex panel use reduces time to diagnosis, duration of antibiotic therapy, length of hospital stay, and need for additional diagnostic testing including chest radiography in some settings, though impact on overall mortality has not been consistently demonstrated.

\section*{Why This Test Is Ordered}
\begin{itemize}
  \item Comprehensive, simultaneous detection of viral and bacterial respiratory pathogens in hospitalized patients with community-acquired or healthcare-associated pneumonia, where rapid identification of the pathogen can guide targeted therapy and antibiotic stewardship
  \item Rapid diagnostic evaluation of immunocompromised patients (HSCT, solid organ transplant, neutropenic) with fever and respiratory symptoms, where identification of a viral pathogen may avoid invasive diagnostic procedures and allow early antiviral therapy
  \item Evaluation of patients with severe respiratory illness requiring ICU admission, where rapid multiplex results inform infection control, antibiotic decisions, and prognostication
  \item Pediatric patients with bronchiolitis or pneumonia who require hospitalization, allowing pathogen-specific cohorting and avoidance of unnecessary antibiotics in confirmed viral infection
  \item Investigation of patients with persistent or worsening respiratory illness in whom initial single-plex testing (influenza, RSV, SARS-CoV-2) was negative
  \item Surveillance for respiratory pathogen epidemiology in the community and hospital settings, including identification of emerging pathogens and outbreak detection
  \item Antimicrobial stewardship: positive identification of a viral pathogen can support the decision to withhold or discontinue empiric antibiotics in patients with bronchiolitis or viral pneumonia
\end{itemize}

\section*{Is This a Primary or Secondary Test}
This is a primary diagnostic test for comprehensive respiratory pathogen detection in hospitalized patients with moderate-to-severe respiratory illness. However, due to its higher cost ($200--500+) compared to single-target testing ($30--100), many institutions restrict its use to specific patient populations (ICU, immunocompromised, pediatric) or require approval. Clinical guidelines from the IDSA and ATS do not mandate multiplex panel use for community-acquired pneumonia but suggest it may be useful for epidemiologic purposes and antimicrobial stewardship.

\section*{How This Test Is Performed}
Specimen collection: Nasopharyngeal (NP) flocked swab is the standard sample type for most platforms. The swab is inserted into one nostril to the posterior nasopharynx, rotated for 5--10 seconds, and placed in the manufacturer-provided transport medium (not standard viral transport medium for FilmArray; the specific sample buffer supplied with cartridges must be used). Combined nasopharyngeal and oropharyngeal swabs may improve sensitivity. Nasopharyngeal wash/aspirate is acceptable for most platforms, particularly in pediatric patients. For intubated patients, endotracheal aspirate, mini-BAL, or bronchoalveolar lavage (BAL) specimens can be used. Each platform has specific specimen requirements, volumes (typically 200--350 $\mu$L of specimen in transport medium), and collection devices detailed in the package insert.

\section*{What Preparation Is Needed}
No special patient preparation is required. Avoid use of intranasal sprays or saline immediately before specimen collection. For platforms using their specific transport medium vial (FilmArray, ePlex), the specimen swab must be placed directly into the manufacturer's vial, not into standard universal transport medium.

\section*{Contraindications and Precautions}
Results must be interpreted in the context of the clinical presentation. Detection of a respiratory virus does not exclude bacterial co-infection, which occurs in 10--30\% of patients with viral pneumonia. Detection of certain pathogens (e.g., rhinovirus/enterovirus, adenovirus) may represent asymptomatic shedding or prolonged detection after resolution of acute infection rather than the cause of current illness. The FilmArray RP detects adenovirus with reduced sensitivity for serotypes in species D and F. False-negative results may occur with low organism load, mutations in primer/probe binding sites, or specimens collected late in illness. The panel does not include all respiratory pathogens (e.g., \textit{Legionella} species, non-respiratory coronaviruses) and does not replace culture for susceptibility testing.

\section*{Risks}
Nasopharyngeal specimen collection may cause transient discomfort, sneezing, coughing, gagging, or mild epistaxis. When collecting specimens from patients with suspected highly transmissible respiratory pathogens (SARS-CoV-2, measles), use of N95 respirator and eye protection is required given the risk of aerosolization.

\section*{What Happens After the Test}
Results are available in 45 minutes (BioFire FilmArray RP2.1), approximately 1 hour (QIAstat-Dx), approximately 2 hours (GenMark ePlex), or approximately 4 hours (Luminex NxTAG). For FilmArray, the pouch can be loaded and started at the point of care by non-laboratory personnel (CLIA-waived status). Interpretation requires correlation with clinical findings; detection of multiple pathogens may indicate co-infection, sequential infection, or carriage of one pathogen without active disease.

\section*{Understanding Results}

\subsection*{Below Normal}
\begin{itemize}
  \item \textbf{Panel negative (all targets not detected):} Indicates that none of the panel pathogens was detected in the specimen. This may represent: non-infectious cause of respiratory symptoms (aspiration pneumonitis, pulmonary edema, pulmonary embolism, drug-induced pneumonitis); infection with a pathogen not included in the panel (\textit{Legionella pneumophila}, \textit{Pneumocystis jirovecii}, \textit{Mycobacterium tuberculosis}, endemic fungi, Nocardia); bacterial pneumonia with a low organism load in the upper respiratory tract; or specimen collected after viral clearance
  \item Up to 40--60\% of community-acquired pneumonia cases remain undiagnosed even with comprehensive multiplex testing, reflecting the limitations of upper respiratory sampling and the presence of pathogens not included in current panels
\end{itemize}

\subsection*{Above Normal}
\begin{itemize}
  \item \textbf{Single pathogen detected:} Most straightforward interpretation. Identification of influenza A/B, RSV, parainfluenza, human metapneumovirus, adenovirus, rhinovirus, or SARS-CoV-2 in a patient with compatible illness strongly supports the causative diagnosis. Identification of \textit{B. pertussis}, \textit{M. pneumoniae}, or \textit{C. pneumoniae} supports atypical bacterial etiology requiring macrolide or tetracycline therapy
  \item \textbf{Multiple pathogens detected:} Co-detection of 2 or more viruses occurs in 10--20\% of positive panels, most commonly in young children. The clinical significance of co-detection is uncertain; some may represent true co-infections, others may reflect prolonged shedding of a prior infection alongside a new acute pathogen. Rhinovirus/enterovirus is the most frequent co-detected virus
  \item \textbf{Rhinovirus/enterovirus detected:} These are the most commonly detected respiratory viruses but also the most challenging to interpret. Rhinovirus is detected in 10--40\% of asymptomatic children and can be shed for 2--3 weeks after resolution of symptoms. Detection should be interpreted cautiously as the sole cause of severe lower respiratory illness
  \item \textbf{Bacterial targets detected:} Detection of \textit{B. pertussis}, \textit{B. parapertussis}, \textit{M. pneumoniae}, or \textit{C. pneumoniae} has high positive predictive value for infection, as asymptomatic carriage is uncommon
\end{itemize}

\section*{Normal Results}
\textbf{Negative} or \textbf{Not detected} for all panel targets. No respiratory pathogen nucleic acid detected.

\section*{Follow-up Tests}
\begin{itemize}
  \item \textbf{Blood cultures and sputum culture:} When the panel is negative but bacterial pneumonia is clinically suspected; cultures also provide antimicrobial susceptibility data that molecular panels do not
  \item \textbf{Legionella urinary antigen:} When \textit{Legionella} is suspected (severe CAP, hyponatremia, elevated liver enzymes, recent travel/hotel stay) as it is not on most multiplex panels
  \item \textbf{Procalcitonin:} Adjunctive biomarker that may help distinguish bacterial from viral infection when multiplex panel results are ambiguous; procalcitonin <0.25 ng/mL supports a viral etiology and may justify withholding antibiotics
  \item \textbf{Single-target confirmatory PCR:} When a high-consequence pathogen (SARS-CoV-2) or public health notifiable disease is detected, confirmatory testing may be required per state and institutional protocols
  \item \textbf{Legionella culture, fungal culture, AFB culture:} When multiplex panel is negative in immunocompromised patients with progressive pneumonia
\end{itemize}

\section*{References}
\begin{enumerate}
  \item Hanson KE, Azar MM, Banerjee R, et al. Molecular testing for acute respiratory tract infections: clinical and diagnostic recommendations from the IDSA and AMP. Clin Infect Dis. 2020;71(10):2744-2751.
  \item Leber AL, Lisby JG, Hansen G, et al. Multicenter evaluation of the BioFire FilmArray Respiratory Panel 2 for detection of viruses and bacteria in nasopharyngeal swab samples. J Clin Microbiol. 2018;56(6):e01945-17.
  \item Rogers BB, Shankar P, Jerris RC, et al. Impact of a rapid respiratory panel test on patient outcomes. Arch Pathol Lab Med. 2015;139(5):636-641.
  \item Brendish NJ, Malachira AK, Armstrong L, et al. Routine molecular point-of-care testing for respiratory viruses in adults presenting to hospital with acute respiratory illness (ResPOC): a pragmatic, open-label, RCT. Lancet Respir Med. 2017;5(5):401-411.
  \item Uyeki TM, Bernstein HH, Bradley JS, et al. Clinical practice guidelines by the IDSA: 2018 update on diagnosis, treatment, chemoprophylaxis of seasonal influenza. Clin Infect Dis. 2019;68(6):e1-e47.
\end{enumerate}

\clearpage
